\documentclass[11pt]{article}

% ---------- Packages ----------
\usepackage[utf8]{inputenc}
\usepackage[T1]{fontenc}
\usepackage{lmodern}
\usepackage[margin=1in]{geometry}
\usepackage{amsmath,amssymb}
\usepackage{booktabs}
\usepackage{array}
\usepackage{siunitx}
\usepackage{graphicx}
\usepackage{caption}
\usepackage[dvipsnames]{xcolor}
\usepackage{titlesec}
\usepackage[hidelinks]{hyperref}

% Absorb slack from long monospace tokens in justified text (avoids overfull boxes)
\emergencystretch=3em
\hbadness=10000

% ---------- Style ----------
\definecolor{headingblue}{RGB}{31,73,125}
\titleformat{\section}{\normalfont\Large\bfseries\color{headingblue}}{\thesection.}{0.5em}{}
\titleformat{\subsection}{\normalfont\large\bfseries\itshape\color{headingblue}}{\thesubsection}{0.5em}{}
\setlength{\parskip}{0.4em}
\setlength{\parindent}{1.2em}
\captionsetup{font=small,labelfont={bf,it},textfont=it}
\sisetup{round-mode=places,round-precision=3,table-format=3.3}

% ---------- Title ----------
\title{\vspace{-1.5em}\bfseries\color{headingblue}\Large
Predicting Superconducting Critical Temperature from Engineered
Material Descriptors: A Comparative Machine Learning Study}

\author{
Israfil Saygitov\thanks{Independent Researcher. Corresponding author.}
\and
Nursan Omarov\thanks{Independent Researcher.}
}
\date{}

\begin{document}
\maketitle

% ---------- Abstract ----------
\begin{abstract}
\noindent
Superconducting materials, which conduct electric current with zero resistance below a critical
temperature ($T_c$), represent one of the most promising avenues for reducing energy losses in
electrical infrastructure. Identifying material compositions associated with a higher $T_c$ remains a
central open problem in condensed matter physics and materials science. This study investigates whether
machine learning models can predict $T_c$ from a set of $81$ engineered numerical descriptors derived
from elemental and compositional properties. Six regression models --- Linear Regression, Decision Tree,
Random Forest, Gradient Boosting, XGBoost, and a deep multilayer perceptron (MLP) --- were trained and
evaluated on a public superconductivity dataset ($21{,}263$ samples after cleaning) using an 80/20
train--test split and 5-fold cross-validation. Random Forest achieved the strongest performance
(cross-validated RMSE $\approx 9.70$~K, test $R^2 \approx 0.930$), narrowly ahead of XGBoost
(CV-RMSE $\approx 9.88$~K), while the deep MLP placed in the middle of the field
(test $R^2 \approx 0.878$) and Linear Regression performed weakest ($R^2 \approx 0.738$), indicating a
strongly nonlinear relationship between descriptors and $T_c$ on which tree ensembles remain especially
strong. Impurity-based and permutation feature importance agree that thermal-conductivity-,
atomic-mass-, and valence-related descriptors are the most informative predictors, and Principal
Component Analysis confirms that the dataset contains non-random structure related to $T_c$. We
additionally address a data-leakage risk: the $21{,}263$ measurements span only $15{,}542$ distinct chemical
formulas, so a random split shares compositions between train and test. Under a rigorous composition-grouped
split (no chemical formula shared across partitions), Random Forest performance declines only modestly, from
$R^2 \approx 0.930$ to $R^2 \approx 0.913$, and the model ranking is preserved --- indicating that the
models generalize to genuinely unseen compositions. We discuss the limitations of this proof-of-concept
pipeline and outline directions for developing it into a screening tool for candidate superconducting
materials.

\vspace{0.6em}
\noindent\textit{Keywords:} superconductivity; critical temperature; machine learning; random forest;
XGBoost; deep learning; feature importance; materials informatics.
\end{abstract}

% ---------- 1. Introduction ----------
\section{Introduction}

\subsection{Motivation}
In 2026, severe flooding struck the Republic of Dagestan, Russia, and local officials attributed part of
the resulting infrastructure damage to degraded electrical wiring --- cables that wear out, overheat, and
fail under stress. This observation motivated the present study: what if electrical grids did not lose
energy to resistance at all?

Superconductors are materials that conduct electric current with zero electrical resistance, meaning they
do not dissipate energy as heat and, in principle, do not degrade through resistive overheating in the way
conventional conductors do. Widespread use of superconducting materials in power infrastructure could
substantially improve grid reliability and efficiency. The central obstacle to this vision is that most
known superconductors only exhibit this property at extremely low temperatures, which makes them
impractical for large-scale real-world deployment.

Identifying materials with a higher critical temperature ($T_c$) --- the temperature below which
superconductivity occurs --- is therefore one of the most important open problems in modern condensed
matter physics and materials science. This work applies machine learning methods to predict $T_c$ from
engineered material descriptors, as an initial step toward a computational screening pipeline for
promising superconducting candidates.

\subsection{Objectives}
The specific objectives of this study are to: (i) compare the predictive performance of six regression
algorithms --- including a deep neural network --- on a public superconductivity dataset; (ii) identify
which engineered descriptors are most informative for predicting $T_c$; and (iii) assess, through
correlation analysis, dimensionality reduction, and a duplicate-leakage check, whether the dataset
exhibits structure consistent with a learnable physical relationship and how robust the reported
performance is.

% ---------- 2. Data and Methods ----------
\section{Data and Methods}

\subsection{Dataset}
The dataset used in this study consists of engineered numerical descriptors for superconducting materials,
together with the corresponding critical temperature (\texttt{critical\_temp}) as the target variable. Each
descriptor is derived from elemental and compositional properties of the constituent atoms (e.g., atomic
mass, valence, thermal conductivity, electron affinity, atomic radius, and fusion heat), aggregated using
statistics such as mean, weighted mean, geometric mean, standard deviation, range, and entropy across the
elements present in each material. The raw dataset contained $21{,}263$ usable records after
preprocessing, with $81$ descriptor features. The data originate from the UCI Machine Learning Repository
superconductivity dataset of Hamidieh~\cite{hamidieh2018}.

\subsection{Preprocessing}
One column, \texttt{wtd\_std\_Density}, was coerced to a numeric type, and any rows with resulting invalid
or missing values were removed. The target variable was set to \texttt{critical\_temp}, and all remaining
numerical columns were used as input features. The cleaned dataset was split into training ($80\%$) and
test ($20\%$) subsets using a fixed random seed to ensure reproducibility.

\subsection{Models}
Six regression models were trained and compared:
\begin{itemize}
  \item \textbf{Linear Regression}, as a baseline linear model.
  \item \textbf{Decision Tree Regressor}, a single nonlinear tree-based model.
  \item \textbf{Random Forest Regressor} ($300$ estimators), an ensemble of bagged decision trees.
  \item \textbf{Gradient Boosting Regressor}, a sequential boosting ensemble.
  \item \textbf{XGBoost Regressor} ($300$ estimators, max depth $6$, learning rate $0.05$,
        \texttt{subsample} $0.8$, \texttt{colsample\_bytree} $0.8$), a regularized gradient boosting
        implementation.
  \item \textbf{Multilayer Perceptron (MLP)}, a deep feed-forward neural network implemented in PyTorch,
        with three hidden layers ($256$--$128$--$64$ units), batch normalization, ReLU activations, and
        dropout regularization. Inputs and target were standardized using training-set statistics, and the
        network was trained with the Adam optimizer under mean-squared-error loss with early stopping on an
        internal validation split.
\end{itemize}

All models were evaluated on the held-out test set using mean absolute error (MAE), mean squared error
(MSE), root mean squared error (RMSE), and the coefficient of determination ($R^2$). To assess the
stability of each model's performance, 5-fold cross-validation was additionally performed on the training
set, reporting the mean and standard deviation of the cross-validated RMSE (CV-RMSE, CV-STD).

\subsection{Feature Importance and Structural Analysis}
To interpret the best-performing model, impurity-based feature importances were extracted from the Random
Forest model. Because impurity-based importances can be biased when features are correlated, permutation
importance was additionally computed on the held-out test set as a robustness check. In parallel, Pearson
correlation coefficients between each descriptor and \texttt{critical\_temp} were computed across the full
cleaned dataset. Finally, Principal Component Analysis (PCA) was applied to the standardized feature matrix
to visualize the two-dimensional structure of the dataset with respect to critical temperature.

\subsection{Duplicate and Leakage Analysis}
Engineered superconductivity descriptors are known to contain many near-identical compositions. Because a
random train--test split can place near-duplicate materials in both partitions and leak information, we
performed two robustness checks. First, we counted exact-duplicate feature rows and re-evaluated the best
model on a deduplicated dataset. Second, and more rigorously, we obtained the chemical formula of each
material from the companion \texttt{unique\_m.csv} file (which aligns one-to-one with the descriptor rows)
and performed a \emph{composition-grouped} split: all measurements sharing a chemical formula were kept
entirely within either the training or the test partition, so that no formula appears in both. All six
models were re-trained and re-evaluated under this leakage-free protocol using \texttt{GroupShuffleSplit}
for the held-out test set and \texttt{GroupKFold} for 5-fold cross-validation, both keyed on the chemical
formula.

% ---------- 3. Results ----------
\section{Results}

\subsection{Model Comparison}
Table~\ref{tab:models} summarizes the performance of all six models on the test set and under 5-fold
cross-validation. Random Forest achieved the best overall performance (test MAE $=5.08$~K,
RMSE $=8.97$~K, $R^2 = 0.930$; CV-RMSE $=9.70 \pm 0.25$~K), closely followed by XGBoost
(CV-RMSE $=9.88 \pm 0.28$~K); the two tree ensembles are statistically very close. The deep MLP placed in
the middle of the field (test $R^2 = 0.878$, CV-RMSE $=12.11 \pm 0.20$~K), outperforming a single Decision
Tree and Gradient Boosting on cross-validated RMSE but not matching the leading tree ensembles. Linear
Regression performed markedly worse than all nonlinear models ($R^2 = 0.738$, CV-RMSE $=17.70$~K).

\begin{table}[htbp]
\centering
\caption{Performance comparison of six regression models on the test set, sorted by cross-validated RMSE.}
\label{tab:models}
\begin{tabular}{l S S S S S S}
\toprule
\textbf{Model} & {\textbf{MAE}} & {\textbf{MSE}} & {\textbf{RMSE}} & {\textbf{$R^2$}}
 & {\textbf{CV-RMSE}} & {\textbf{CV-STD}} \\
\midrule
Random Forest        & 5.076 & 80.452  & 8.970  & 0.930 & 9.700  & 0.253 \\
XGBoost              & 5.811 & 86.037  & 9.276  & 0.925 & 9.876  & 0.280 \\
MLP (Deep Learning)  & 7.795 & 140.725 & 11.863 & 0.878 & 12.114 & 0.201 \\
Decision Tree        & 6.090 & 139.449 & 11.809 & 0.879 & 12.539 & 0.296 \\
Gradient Boosting    & 8.501 & 152.150 & 12.335 & 0.868 & 12.794 & 0.149 \\
Linear Regression    & 13.211& 302.008 & 17.378 & 0.738 & 17.699 & 0.253 \\
\bottomrule
\end{tabular}
\end{table}

The gap between Linear Regression and the tree-based ensembles suggests that the relationship between the
engineered descriptors and critical temperature is strongly nonlinear. Notably, the deep MLP --- despite
its capacity to model nonlinear interactions --- did not surpass Random Forest or XGBoost, consistent with
the broader observation that gradient-boosted and bagged tree ensembles remain strong baselines for
tabular data with heterogeneous engineered features.

\subsection{Feature Importance}
Table~\ref{tab:impurity} lists the descriptors with the highest impurity-based importance in the Random
Forest model. A single feature, \texttt{range\_ThermalConductivity}, dominates the importance ranking
($0.536$), followed by \texttt{wtd\_gmean\_ThermalConductivity} ($0.126$). The remaining top features
relate to atomic mass, valence, and density, together accounting for a much smaller share of total
importance.

\begin{table}[htbp]
\centering
\caption{Top 8 features by Random Forest impurity-based importance.}
\label{tab:impurity}
\begin{tabular}{l S[table-format=1.3]}
\toprule
\textbf{Feature} & {\textbf{Importance}} \\
\midrule
range\_ThermalConductivity        & 0.536 \\
wtd\_gmean\_ThermalConductivity   & 0.126 \\
std\_atomic\_mass                 & 0.022 \\
wtd\_gmean\_Valence               & 0.019 \\
std\_Density                      & 0.011 \\
mean\_Density                     & 0.011 \\
wtd\_entropy\_ThermalConductivity & 0.011 \\
gmean\_ElectronAffinity           & 0.011 \\
\bottomrule
\end{tabular}
\end{table}

We note that this pronounced concentration of importance in thermal-conductivity-related descriptors may
partly reflect multicollinearity among related features (e.g., \texttt{range\_}, \texttt{wtd\_gmean\_}, and
\texttt{wtd\_std\_ThermalConductivity} are derived from the same underlying property) rather than a single
dominant physical mechanism; this is discussed further in Section~\ref{sec:discussion}.

\subsection{Permutation Importance}
To guard against the known bias of impurity-based importance under feature correlation, permutation
importance was computed on the held-out test set (Table~\ref{tab:perm}). The two measures agree closely:
\texttt{range\_ThermalConductivity} remains dominant ($0.722 \pm 0.019$), followed by
\texttt{wtd\_gmean\_ThermalConductivity}, \texttt{std\_atomic\_mass}, and \texttt{wtd\_gmean\_Valence}. The
agreement between two independent importance measures strengthens confidence that these descriptors carry
genuine predictive signal rather than being artifacts of a single scoring method.

\begin{table}[htbp]
\centering
\caption{Top 8 features by Random Forest permutation importance (mean $\pm$ std over 5 repeats,
held-out test set).}
\label{tab:perm}
\begin{tabular}{l c}
\toprule
\textbf{Feature} & \textbf{Permutation importance} \\
\midrule
range\_ThermalConductivity        & $0.722 \pm 0.019$ \\
wtd\_gmean\_ThermalConductivity   & $0.079 \pm 0.005$ \\
std\_atomic\_mass                 & $0.039 \pm 0.002$ \\
wtd\_gmean\_Valence               & $0.033 \pm 0.001$ \\
mean\_Density                     & $0.018 \pm 0.001$ \\
entropy\_FusionHeat               & $0.015 \pm 0.001$ \\
entropy\_Density                  & $0.014 \pm 0.000$ \\
range\_atomic\_radius             & $0.014 \pm 0.001$ \\
\bottomrule
\end{tabular}
\end{table}

\subsection{Correlation Analysis}
Pearson correlation coefficients between individual descriptors and \texttt{critical\_temp} ranged from
approximately $-0.63$ to $+0.72$. The strongest positive correlations were observed for
\texttt{wtd\_std\_ThermalConductivity} ($r=0.721$), \texttt{range\_ThermalConductivity} ($r=0.688$),
\texttt{range\_atomic\_radius} ($r=0.654$), and \texttt{std\_ThermalConductivity} ($r=0.654$). The
strongest negative correlations were observed for \texttt{wtd\_mean\_Valence} ($r=-0.632$) and
\texttt{wtd\_gmean\_Valence} ($r=-0.616$). These results are broadly consistent with the Random Forest
feature importance ranking and indicate that thermal-conductivity-, radius-, and valence-related
descriptors carry meaningful information about critical temperature, although correlation alone does not
establish causal physical relationships.

\subsection{Distribution and Structure of the Target Variable}
The distribution of \texttt{critical\_temp} across the dataset is right-skewed (skewness $\approx 0.86$;
mean $\approx 34.4$~K, median $\approx 20.0$~K, maximum $185$~K): most materials exhibit relatively low
critical temperatures, with a smaller number of materials reaching substantially higher values. This
asymmetry means the RMSE is dominated by the harder-to-predict high-$T_c$ tail and suggests that the
dataset may contain multiple underlying material regimes with distinct superconducting behavior.

A two-dimensional PCA projection of the standardized feature matrix explained approximately $49.4\%$ of
total variance (PC1: $38.9\%$, PC2: $10.5\%$). Materials with higher critical temperatures tended to
occupy specific regions of this projection rather than being uniformly distributed, indicating that the
engineered descriptors capture structure related to critical temperature, although the moderate explained
variance implies this two-dimensional view captures only part of the dataset's complexity.

\subsection{Duplicate and Leakage Robustness}
A large fraction of the dataset consists of exact-duplicate feature vectors: $6{,}093$ of $21{,}263$ rows
($28.7\%$) are exact duplicates of another row in feature space. Removing these duplicates leaves $15{,}170$
unique rows. Retraining Random Forest on this deduplicated dataset lowers performance from
$R^2 = 0.930$ (RMSE $8.97$~K) to $R^2 = 0.911$ (RMSE $10.01$~K).

The stronger check is the composition-grouped split. The $21{,}263$ measurements correspond to only
$15{,}542$ distinct chemical formulas --- some formulas (e.g., \texttt{Y1Ba2Cu3O7}) appear over one hundred
times --- so a random split shares many formulas between train and test. Under the leakage-free
composition-grouped split (no formula shared across partitions; $4{,}274$ test measurements),
performance drops only modestly, confirming that the models genuinely generalize to unseen compositions
rather than merely memorizing repeated ones. Table~\ref{tab:grouped} reports all six models under both
protocols. Random Forest, the best model, declines from $R^2 = 0.930$ to $R^2 = 0.913$
(RMSE $8.97 \to 10.08$~K, grouped CV-RMSE $9.94$~K), and the overall model ranking is preserved. We regard
the composition-grouped figures as the honest, leakage-controlled estimates of real-world screening
performance.

\begin{table}[htbp]
\centering
\caption{Random (leaky) versus composition-grouped (leakage-free) evaluation. Grouped columns use a
formula-keyed \texttt{GroupShuffleSplit} test set and \texttt{GroupKFold} cross-validation; models sorted by
grouped CV-RMSE.}
\label{tab:grouped}
\begin{tabular}{l S[table-format=1.3] S[table-format=1.3] S[table-format=2.3] S[table-format=2.3] S[table-format=2.3]}
\toprule
 & \multicolumn{2}{c}{\textbf{Test $R^2$}} & \multicolumn{2}{c}{\textbf{Test RMSE (K)}}
 & {\textbf{Grouped}} \\
\cmidrule(lr){2-3}\cmidrule(lr){4-5}
\textbf{Model} & {\textbf{Random}} & {\textbf{Grouped}} & {\textbf{Random}} & {\textbf{Grouped}}
 & {\textbf{CV-RMSE}} \\
\midrule
Random Forest        & 0.930 & 0.913 & 8.970  & 10.077 & 9.938  \\
XGBoost              & 0.925 & 0.909 & 9.276  & 10.321 & 10.001 \\
MLP (Deep Learning)  & 0.878 & 0.858 & 11.863 & 12.885 & 12.232 \\
Gradient Boosting    & 0.868 & 0.854 & 12.335 & 13.049 & 12.796 \\
Decision Tree        & 0.879 & 0.853 & 11.809 & 13.099 & 13.376 \\
Linear Regression    & 0.738 & 0.726 & 17.378 & 17.877 & 17.639 \\
\bottomrule
\end{tabular}
\end{table}

Because this check removes leakage at the level of the chemical formula, the grouped figures should be
regarded as the primary, honest estimate; the small size of the drop is itself an encouraging result.

% ---------- 4. Discussion ----------
\section{Discussion and Limitations}
\label{sec:discussion}
The results indicate that engineered elemental descriptors carry substantial predictive signal for
critical temperature, and that tree-based ensemble methods --- particularly Random Forest and XGBoost ---
capture this signal considerably better than a linear model and, on this dataset, better than a
conventional deep neural network. However, several limitations should be considered when interpreting
these results.

\begin{itemize}
  \item The dataset uses engineered statistical descriptors rather than direct experimental or
        first-principles physical measurements, which limits the physical interpretability of the learned
        relationships.
  \item A large fraction of rows are exact-duplicate feature vectors ($28.7\%$), and the $21{,}263$
        measurements span only $15{,}542$ distinct chemical formulas, so a random split leaks similar
        compositions across train and test. We address this directly with a composition-grouped split
        (Section~3.6); the residual $R^2$ drop is small, but grouped performance ($R^2 \approx 0.913$ for
        Random Forest) should be read as the realistic estimate rather than the random-split figure.
  \item Feature importance and correlation analysis identify statistically predictive descriptors but do
        not establish causal physical mechanisms, particularly given the strong multicollinearity among
        thermal-conductivity-related features.
  \item Hyperparameters for the boosting methods and the neural network were only lightly tuned, so the
        relative ranking among the top models may shift under more systematic tuning.
  \item PCA is a linear dimensionality-reduction technique used here purely for visualization; it does not
        provide a direct physical interpretation of the principal components.
\end{itemize}

For these reasons, the present results should be regarded as a proof of concept rather than a definitive
scientific conclusion about the physical determinants of superconducting critical temperature.

% ---------- 5. Future Work ----------
\section{Future Work}
\begin{itemize}
  \item Testing the pipeline on additional, independently sourced superconductivity datasets.
  \item Using composition-grouped or deduplicated splitting to obtain leakage-free performance estimates.
  \item Combining engineered descriptors with more direct physical or first-principles measurements
        (e.g., density-functional-theory-derived features).
  \item Systematic hyperparameter tuning for all models, including the neural network architecture.
  \item Developing a ranking-based screening pipeline to prioritize candidate materials for experimental
        investigation.
  \item Adding uncertainty estimation (e.g., quantile regression or ensemble variance) to model
        predictions.
\end{itemize}

% ---------- 6. Conclusion ----------
\section{Conclusion}
This study demonstrates that machine learning models can predict superconducting critical temperature from
engineered material descriptors with meaningful accuracy. Among six compared models, Random Forest achieved
the best performance (cross-validated RMSE $\approx 9.70$~K, $R^2 \approx 0.930$), narrowly outperforming
XGBoost, while a deep multilayer perceptron placed in the middle of the field --- confirming that
gradient-boosted and bagged tree ensembles remain strong baselines for tabular materials data. The
magnitude of improvement over the linear baseline indicates that the relationship between engineered
descriptors and critical temperature is strongly nonlinear. Impurity-based and permutation feature
importance, correlation analysis, and PCA visualization consistently point to thermal-conductivity-,
valence-, mass-, and radius-related descriptors as informative predictors, though these findings should be
interpreted as statistical associations rather than confirmed physical mechanisms. Crucially, a
composition-grouped evaluation --- in which no chemical formula is shared between training and test ---
shows that the models generalize to genuinely unseen compositions, with Random Forest declining only from
$R^2 \approx 0.930$ to $R^2 \approx 0.913$; we regard this leakage-controlled figure as the honest estimate
of real-world screening performance. Overall,
this work serves as a proof-of-concept pipeline for machine-learning-assisted superconductor screening and
provides a foundation for future extensions toward more physically grounded materials discovery workflows.

% ---------- References ----------
\begin{thebibliography}{9}
\bibitem{hamidieh2018}
Hamidieh, K. A data-driven statistical model for predicting the critical temperature of a superconductor.
\textit{Computational Materials Science}, 154, 346--354 (2018).

\bibitem{scikit}
Pedregosa, F. et al. Scikit-learn: Machine Learning in Python.
\textit{Journal of Machine Learning Research}, 12, 2825--2830 (2011).

\bibitem{xgboost}
Chen, T. \& Guestrin, C. XGBoost: A Scalable Tree Boosting System.
\textit{Proceedings of the 22nd ACM SIGKDD International Conference on Knowledge Discovery and Data Mining},
785--794 (2016).

\bibitem{rf}
Breiman, L. Random Forests. \textit{Machine Learning}, 45(1), 5--32 (2001).

\bibitem{pytorch}
Paszke, A. et al. PyTorch: An Imperative Style, High-Performance Deep Learning Library.
\textit{Advances in Neural Information Processing Systems}, 32, 8024--8035 (2019).

\bibitem{pca}
Jolliffe, I. T. \& Cadima, J. Principal component analysis: a review and recent developments.
\textit{Philosophical Transactions of the Royal Society A}, 374(2065) (2016).
\end{thebibliography}

\end{document}
